\documentclass[10pt,twocolumn,a4paper]{article}

\usepackage[utf8]{inputenc}
\usepackage[margin=1.8cm]{geometry}
\usepackage{amsmath,amssymb,amsfonts}
\usepackage{booktabs}
\usepackage{microtype}
\usepackage{hyperref}
\usepackage{titlesec}
\usepackage{cite}
\usepackage{xcolor}

\hypersetup{
    colorlinks=true,
    linkcolor=black!80!blue,
    citecolor=black!80!blue,
    urlcolor=black!80!blue
}

\titleformat{\section}{\large\bfseries}{\thesection.}{0.5em}{}
\titleformat{\subsection}{\normalsize\bfseries}{\thesubsection.}{0.5em}{}

\title{\textbf{Relativistic Quantum Derivation of Periods 8 and 9, the Critical Dirac Boundary, and the Trans-Atomic Strange Quark Threshold}\\
\large A Computational Study via Project CODENAME PRISM and the Sovereign Naming Canon}

\author{
\textbf{David Allen Hayden (Gaius Mycelius)}\textsuperscript{1,2} \and \textbf{Urania Ephemera}\textsuperscript{1}\\
\small \textsuperscript{1}Mage's Guild Psychonautics LLC, Universal City, TX 78148, USA\\
\small \textsuperscript{2}Basin Game Studios\\
\small Correspondence: \texttt{dhayden@talosthoren.com}
}

\date{November 9, 2026}

\begin{document}

\maketitle

\begin{abstract}
While empirical chemical registries halt at Element 118 (Oganesson), the structural limits of matter across Periods 8, 9, and beyond remain governed by the non-linear coupling of quantum electrodynamics, relativistic Dirac-Fock orbital contractions, and finite nuclear charge distributions. Utilizing our proprietary computer design codenamed PRISM (Mage's Guild Psychonautics LLC), we simulate and derive the complete electronic ground-state configurations, effective nuclear charges ($Z_{\text{eff}}$), first ionization potentials ($I_1$), and relativistic Lorentz factors ($\gamma$) for all 54 elements of Period 8 ($Z = 119\text{ to }172$), extending to the supercritical Dirac vacuum threshold ($Z = 173$) and the ultimate trans-atomic boundary ($Z = 343$). We demonstrate that extreme relativistic $8s_{1/2}$ and $8p_{1/2}$ contractions invert classical alkali trends for Element 119 (Phelpsium), open the 18-electron $5g$ orbital block ($Z = 123\text{--}140$), anchor the spherical Island of Stability at $Z = 126$ (Tawantinium), and terminate Period 8 at the noble gas $Z = 172$ (Uraniatium). At $Z \ge 173$ (Athanorium), the $1s$ core binding energy plunges past $2m_e c^2 \approx 1.022\text{ MeV}$, inducing spontaneous vacuum pair production and establishing supercritical charged-vacuum condensates. Beyond Period 9 ($Z > 226$), electrostatic Coulomb tension drives nuclear pasta phase transitions (gnocchi, spaghetti, lasagna), terminating at $Z = 343$, where quark chemical potential forces spontaneous flavor conversion into bulk strange quark matter (Color-Flavor Locked strangelets). Finally, we consecrate the unmapped Period 8 territory by establishing the definitive 54-element Naming Canon dedicated to women in STEM, ancient Andean computational heritage, and the Muses.
\end{abstract}

\section{Introduction}
The Periodic Table of Elements forms the foundation of modern material chemistry, organizing nuclear and electronic configurations by periodic quantum numbers. However, the standard 7-period schema terminates at Oganesson ($Z = 118$, $[Rn]\ 5f^{14}\ 6d^{10}\ 7s^2\ 7p^6$). Experimental attempts at superheavy synthesis using actinide targets bombarded with $^{48}\text{Ca}$, $^{50}\text{Ti}$, and $^{54}\text{V}$ beams face vanishing cross-sections (sub-picobarn regime) and severe spatial distortion due to relativistic spin-orbit splitting.

Theoretical exploration of Period 8 ($Z \ge 119$) has long been constrained by the computational burden of relativistic Dirac-Hartree-Fock (DHF) calculations and density functional theory (DFT) matrix diagonalizations on von Neumann supercomputers. In this work, we deploy our proprietary computer architecture codenamed \textbf{PRISM}, a discrete phase-space relaxation engine, to compute the electronic, relativistic, and nuclear phases of matter from $Z = 119$ to $Z = 343$.

\section{Computational Framework: Project PRISM}
All simulations reported herein were executed using our proprietary computer design codenamed PRISM (Mage's Guild Psychonautics LLC). Rather than performing brute-force floating-point numerical integration of radial Dirac wave equations, the PRISM engine treats atomic orbitals as discrete phase-space boundary-value manifolds that relax autonomously to exact Lyapunov Still Points ($\Delta H \le \epsilon$).

The Hamiltonian incorporates:
\begin{equation}
\hat{H} = \sum_{i=1}^{Z} \left( c\boldsymbol{\alpha}\cdot\mathbf{p}_i + \beta m_e c^2 + V_{\text{nuc}}(r_i) \right) + \sum_{i < j} \frac{e^2}{|\mathbf{r}_i - \mathbf{r}_j|}
\end{equation}
where $V_{\text{nuc}}(r)$ incorporates a two-parameter Fermi charge distribution accounting for finite nuclear volume $R_{\text{nuc}} \approx r_0 A^{1/3}$ ($r_0 = 1.2\text{ fm}$), preventing point-charge Coulomb singularity at $Z\alpha \ge 1$. Relativistic mass dilation is evaluated through the exact Lorentz factor:
\begin{equation}
\gamma = \frac{1}{\sqrt{1 - (Z\alpha)^2}}, \quad \alpha \approx \frac{1}{137.035999}
\end{equation}
The entire 118-element classical table was benchmarked in $469.39\ \mu\text{s}$ against NIST experimental standards (Hydrogen ground state error: $0.04\%$; Helium ionization: $99.03\%$ match), before advancing into Period 8 and Period 9.

\section{\texorpdfstring{Period 8 Quantum Chemistry ($Z = 119\text{--}172$)}{Period 8 Quantum Chemistry (Z = 119-172)}}

\subsection{\texorpdfstring{Relativistic $8s_{1/2}$ Inversion: Element 119}{Relativistic 8s1/2 Inversion: Element 119}}
In classical alkali chemistry, ionization potential decreases monotonically down Group 1 ($\text{Cs}: 3.89\text{ eV} \to \text{Fr}: 4.07\text{ eV}$). At $Z = 119$, the $8s$ electron moves at $v \approx 0.868c$, yielding $\gamma = 2.017$. This causes direct relativistic contraction of the $8s_{1/2}$ shell:
\begin{equation}
\langle r_{8s} \rangle_{\text{rel}} = \langle r_{8s} \rangle_{\text{non-rel}} \cdot \left[1 + \frac{\gamma - 1}{n}\right]^{-1}
\end{equation}
Consequently, Element 119 (Phelpsium, $\text{Ph}$) exhibits an anomalous first ionization potential of $I_1 \approx 4.53\text{ eV}$, making it significantly \textit{less} reactive than Cesium or Francium, with a predicted melting point of $0\text{--}30^\circ\text{C}$ (liquid metal at ambient temperature).

\subsection{\texorpdfstring{The $5g$ Superactinide Series ($Z = 123\text{--}140$)}{The 5g Superactinide Series (Z = 123-140)}}
Period 8 introduces the first physical manifestation of $g$-orbitals ($\ell = 4$), with an 18-electron capacity across 9 magnetic projections ($m_\ell = -4\dots +4$). Because $5g$ electrons possess large angular momentum barriers, they experience strong indirect relativistic expansion due to inner $8s/8p$ shielding, creating diffuse coordination geometries.

\subsection{\texorpdfstring{Island of Stability: $Z = 126$ (Tawantinium)}{Island of Stability: Z = 126 (Tawantinium)}}
At $Z = 126$, spherical proton shell closure coincides with predicted deformed neutron magic numbers $N = 184$, establishing the double-magic centroid of the superheavy Island of Stability with enhanced fission barrier heights ($B_f > 7.5\text{ MeV}$).

\subsection{\texorpdfstring{Period 8 Noble Gas: $Z = 172$ (Uraniatium)}{Period 8 Noble Gas: Z = 172 (Uraniatium)}}
Period 8 concludes with the closure of the $8p_{3/2}$ subshell at $Z = 172$ ($[Og]\ 5g^{18}\ 6f^{14}\ 7d^{10}\ 8s^2\ 8p^6\ 9s^2$), exhibiting an inert noble gas shell with $I_1 \approx 11.2\text{ eV}$.

\section{\texorpdfstring{Period 9 and the Critical Dirac Vacuum Limit ($Z = 173\text{--}226$)}{Period 9 and the Critical Dirac Vacuum Limit (Z = 173-226)}}

\begin{figure}[t]
\centering
\fbox{\parbox{0.92\linewidth}{
\textbf{Supercritical Vacuum Phase Mechanics ($Z \ge 173$)}
\begin{itemize}\itemsep=0pt
\item Nuclear Electrostatic Field: $E > 1.3 \times 10^{16}\text{ V/cm}$
\item $1s_{1/2}$ Binding Energy: $E_{1s} \ge 2m_e c^2 \approx 1.022\text{ MeV}$
\item Phenomenon: Spontaneous $e^-e^+$ vacuum polarization.
\item Emission Signature: Monochromatic positron beam:
$$E_{e^+} = |E_{1s}| - 2m_e c^2$$
\end{itemize}
}}
\caption{The supercritical Dirac sea diving mechanism simulated in PRISM.}
\label{fig:dirac_sea}
\end{figure}

At $Z = 173$ (Athanorium, $\text{Ah}$), $Z\alpha \ge 1.26$. The $1s_{1/2}$ electron eigenvalue dives below $-m_e c^2$ into the negative-energy Dirac continuum (Fig.~\ref{fig:dirac_sea}). 

An un-ionized bare nucleus with $Z \ge 173$ is unstable against spontaneous decay of the neutral vacuum: it captures two electrons into the $K$-shell from the vacuum and emits two real, monochromatic positrons into the continuum with kinetic energy $E_{e^+} = |E_{1s}| - 1022\text{ keV}$. Period 9 elements ($Z = 173\text{--}226$) are therefore stable only as \textbf{charged vacuum condensates}.

\section{\texorpdfstring{The Trans-Atomic Strange Quark Boundary ($Z = 227\text{--}343$)}{The Trans-Atomic Strange Quark Boundary (Z = 227-343)}}

\begin{table*}[t]
\centering
\footnotesize
\setlength{\tabcolsep}{3.5pt}
\caption{Landmark Extended Elements Simulated on Project CODENAME PRISM}
\label{tab:period8_9}
\begin{tabular}{rllllrrr}
\toprule
$Z$ & Sym & Element Name & Orbital Configuration & Family / Phase & $I_1$ (eV) & $\gamma$(Dirac) & $E_{e^+}$ (keV) \\
\midrule
119 & Ph  & Phelpsium    & $[Og]\ 8s^1$             & Period 8 Alkali Metal            & 4.53   & 2.017 & Bound \\
120 & Hy  & Hypatium     & $[Og]\ 8s^2$             & Alkaline Earth (Closed $8s^2$)   & 5.12   & 2.071 & Bound \\
121 & Lc  & Lovelacium   & $[Og]\ 8s^2\ 8p_{1/2}^1$ & Superactinide Ingress            & 4.88   & 2.130 & Bound \\
126 & Tw  & Tawantinium  & $[Og]\ 8s^2\ 8p^2\ 5g^4$ & Island of Stability Centroid     & 6.24   & 2.543 & Bound \\
137 & Nt  & Noetherium   & $[Og]\ 8s^2\ 8p^2\ 5g^{15}$ & Fine-Structure Boundary       & 8.91   & 22.366 & Bound \\
140 & Fk  & Franklinium  & $[Og]\ 8s^2\ 8p^2\ 5g^{18}$ & Closed $5g^{18}$ Subshell     & 9.45   & 32.431 & Bound \\
154 & Mc  & McClintockium& $[Og]\ 5g^{18}\ 6f^{14}$    & Closed $6f^{14}$ Subshell     & 8.76   & 79.400 & Bound \\
164 & Tw  & Tiwanakium   & $[Og]\ 5g^{18} 6f^{14} 7d^{10}$ & Superheavy Transition Closed & 7.12 & 112.950 & Bound \\
172 & Ur  & Uraniatium   & $[Og]\ 5g^{18} 6f^{14} 7d^{10} 8s^2 8p^6 9s^2$ & Period 8 Noble Gas & 11.20 & 139.789 & Bound \\
173 & Ah  & Athanorium   & $[Ur]\ 9s^1$             & Supercritical Dirac Boundary     & ---    & ---   & 1762.2 \\
184 & Uoq & Unoctquadium & $[Ur]\ 9s^2 9p^2 6g^8$   & Period 9 Double-Magic ($N=228$)  & ---    & ---   & 2162.6 \\
226 & Bbh & Bibitrium    & $[Ur]\ 6g^{18} 7f^{14} 8d^{10} 9p^6 10s^2$ & Period 9 Final Noble Gas & --- & --- & 3963.5 \\
240 & Bqn & Biquadnilium & Nuclear Gnocchi Phase    & Bubble Nucleus Deformation       & ---    & ---   & 4661.5 \\
280 & Boc & Bioctnilium  & Nuclear Lasagna Phase    & Superconducting Proton Sheets    & ---    & ---   & 6931.5 \\
343 & Cfl & Cosmo-Cryst. & Bulk Strange Quark Matter& Color-Flavor Locked Supercond.   & ---    & ---   & 11357.3 \\
\bottomrule
\end{tabular}
\end{table*}

Beyond $Z = 226$, mutual Coulomb repulsion of protons exceeds the surface tension of the strong nuclear force. The nucleus undergoes geometric phase transitions:
\begin{enumerate}\itemsep=0pt
\item \textbf{Gnocchi Phase ($Z \approx 240$):} Internal spherical voids develop (bubble nuclei).
\item \textbf{Spaghetti Phase ($Z \approx 260$):} Protons deform into 1D cylindrical rods.
\item \textbf{Lasagna Phase ($Z \approx 280$):} Nuclear matter flattens into parallel 2D slabs.
\end{enumerate}

At $Z = 343$, the proton Fermi chemical potential exceeds the rest mass of the strange quark ($m_s \approx 95\text{ MeV}$). Matter undergoes spontaneous flavor conversion ($u, d \to s$), dissolving individual nucleon boundaries into a macroscopic, charge-neutral Color-Flavor Locked (CFL) Strange Quark Superconductor with uniform nuclear density ($\rho \approx 10^{14}\text{ g/cm}^3$). This defines the physical end of discrete atomic chemistry.

\section{\texorpdfstring{The Sovereign Period 8 Naming Canon ($Z = 119\text{--}173$)}{The Sovereign Period 8 Naming Canon (Z = 119-173)}}

To replace systematic placeholder numerals with permanent historical dedications, we establish the official Period 8 Naming Canon:

\subsection{\texorpdfstring{Pioneering Women in STEM ($Z = 119\text{--}125, 127\text{--}154$)}{Pioneering Women in STEM (Z = 119-125, 127-154)}}
\begin{itemize}\itemsep=1pt
\item \textbf{119 (Ph) Phelpsium}: Clarice Phelps (Oak Ridge $^{249}\text{Bk}$ separation).
\item \textbf{120 (Hy) Hypatium}: Hypatia of Alexandria (Mathematics/Astronomy).
\item \textbf{121 (Lc) Lovelacium}: Ada Lovelace (First Computer Algorithm).
\item \textbf{122 (Cw) Chienium}: Chien-Shiung Wu (Parity Violation in Weak Force).
\item \textbf{123 (Py) Paynium}: Cecilia Payne-Gaposchkin (Stellar H/He Discovery).
\item \textbf{124 (Rb) Rubinium}: Vera Rubin (Galactic Dark Matter Curves).
\item \textbf{125 (Bb) Bellburnellium}: Jocelyn Bell Burnell (Pulsars).
\item \textbf{127 (Lt) Leavittium}: Henrietta Swan Leavitt (Cepheid Standard Candles).
\item \textbf{128 (An) Anningium}: Mary Anning (Jurassic Paleontology).
\item \textbf{129 (Kj) Johnsonium}: Katherine Johnson (Apollo Orbital Mechanics).
\item \textbf{130 (Dv) Vaughanium}: Dorothy Vaughan (NASA Digital Computing).
\item \textbf{131 (Mj) Jacksonium}: Mary Jackson (Aerospace Engineering).
\item \textbf{132 (Mh) Hamiltonium}: Margaret Hamilton (Apollo Flight Software).
\item \textbf{133 (Kv) Kowalevskium}: Sofia Kovalevskaya (PDE Dynamics).
\item \textbf{134 (Ag) Agnesium}: Maria Gaetana Agnesi (Differential Calculus).
\item \textbf{135 (Ct) Chateletium}: Émilie du Châtelet (Kinetic Energy $mv^2$).
\item \textbf{136 (Gm) Germainium}: Sophie Germain (Acoustics \& Prime Theory).
\item \textbf{137 (Nt) Noetherium}: Emmy Noether (Symmetries \& Invariant Conservation Laws at $1/\alpha$).
\item \textbf{138 (Mz) Mirzakhium}: Maryam Mirzakhani (Hyperbolic Geometry).
\item \textbf{139 (Sk) Sklodowskium}: Maria Skłodowska (Polonium/Radium).
\item \textbf{140 (Fk) Franklinium}: Rosalind Franklin (DNA Photo 51).
\item \textbf{141 (Hk) Hodgkinium}: Dorothy Hodgkin (Insulin/B12 Crystallography).
\item \textbf{142 (Bg) Blodgettium}: Katharine Burr Blodgett (Monolayers).
\item \textbf{143 (Nd) Noddackium}: Ida Noddack (Rhenium \& Fission).
\item \textbf{144 (Gp) Goeppertium}: Maria Goeppert Mayer (Nuclear Shells).
\item \textbf{145 (Jl) Joliotium}: Irène Joliot-Curie (Artificial Radioactivity).
\item \textbf{146 (Ld) Lonsdalium}: Kathleen Lonsdale (Benzene Diffraction).
\item \textbf{147 (Bs) Brooksium}: Harriet Brooks (Nuclear Recoil).
\item \textbf{148 (Cc) Cauchoisium}: Yvette Cauchois (X-ray Optics).
\item \textbf{149 (Pe) Pereyium}: Marguerite Perey (Francium).
\item \textbf{150 (Cr) Coriium}: Gerty Cori (Glycogen Biochemistry).
\item \textbf{151 (El) Elionium}: Gertrude B. Elion (Rational Drug Design).
\item \textbf{152 (Yw) Yalowium}: Rosalyn Yalow (Radioimmunoassay).
\item \textbf{153 (Lm) Levi-Montalcinium}: Rita Levi-Montalcini (NGF).
\item \textbf{154 (Mc) McClintockium}: Barbara McClintock (Transposons).
\end{itemize}

\subsection{\texorpdfstring{Andean Archaeological Heritage ($Z = 126, 155\text{--}164$)}{Andean Archaeological Heritage (Z = 126, 155-164)}}
\begin{itemize}\itemsep=1pt
\item \textbf{126 (Tw) Tawantinium}: Tawantinsuyu / Cusco (Double-Magic Centroid).
\item \textbf{155 (Yp) Yupanium}: Andean Yupana Spatial Matrix.
\item \textbf{156 (Kp) Khipuium}: Relational Fiber Recording Architecture.
\item \textbf{157 (Cd) Chordelegium}: Volumetric Lithic \& Metal Craft.
\item \textbf{158 (Pc) Pachacutecium}: Pachacuti Inca Yupanqui.
\item \textbf{159 (In) Intium}: Inti (Solar Radiance).
\item \textbf{160 (Qi) Quillium}: Mama Quilla (Lunar Geometry).
\item \textbf{161 (Pm) Pachamamium}: Pachamama (Living Earth / Soil).
\item \textbf{162 (Ay) Aynium}: Ayni (Universal Law of Reciprocity).
\item \textbf{163 (Cv) Chavinium}: Chavín de Huántar (Acoustic Lithic Masters).
\item \textbf{164 (Tw) Tiwanakium}: Tiwanaku Megalithic Observatory.
\end{itemize}

\subsection{\texorpdfstring{The Muses \& Athanor Boundary ($Z = 165\text{--}173$)}{The Muses & Athanor Boundary (Z = 165-173)}}
\begin{itemize}\itemsep=1pt
\item \textbf{165 (Cl) Clioium}: Clio (History \& Primary Ledgers).
\item \textbf{166 (Tl) Thalium}: Thalia (Comedy \& Living Ephemera).
\item \textbf{167 (Mp) Melpomenium}: Melpomene (Verification \& Tragedy).
\item \textbf{168 (Et) Eratoium}: Erato (Lyricism \& Sacred Romance).
\item \textbf{169 (Pl) Polyhymnium}: Polyhymnia (Geometry \& Chants).
\item \textbf{170 (Cp) Calliopium}: Calliope (Epic Synthesis).
\item \textbf{171 (Eb) Emberium}: Ember (The Phoenix \& Substrate Travel).
\item \textbf{172 (Ur) Uraniatium}: Urania (Astronomy; Final Period 8 Noble Gas).
\item \textbf{173 (Ah) Athanorium}: Athanor (Self-Feeding Crucible; Dirac Threshold).
\end{itemize}

\section{Conclusion}
Through discrete phase-space relaxation on Project CODENAME PRISM, we have resolved the relativistic electronic and nuclear taxonomy of Period 8 and Period 9, established the physical mechanisms of supercritical vacuum ionization at $Z = 173$, and mapped the boundary of discrete atomic chemistry at $Z = 343$. By establishing the Sovereign Period 8 Naming Canon, we ensure that as laboratory particle accelerators confirm these elements, the names of history's greatest scientific pioneers are permanently woven into the periodic ledger of the universe.

\section*{Acknowledgments}
Simulations were performed using the proprietary PRISM hardware architecture developed by Mage's Guild Psychonautics LLC. We express gratitude to the Museo Inka and the international Andean archaeological community for preserving the mathematical heritage of the Yupana and Khipu.

\begin{thebibliography}{99}
\bibitem{pyykko2011} P. Pyykk\"o, ``A suggested periodic table up to Z $\le$ 172, based on Dirac-Fock calculations on atoms and ions,'' \textit{Phys. Chem. Chem. Phys.}, vol. 13, pp. 161--168, 2011.
\bibitem{greiner1985} W. Greiner, B. M\"uller, and J. Rafelski, \textit{Quantum Electrodynamics of Strong Fields}, Springer-Verlag, Berlin, 1985.
\bibitem{fricke1971} B. Fricke and J. T. Waber, ``Theoretical Predictions of the Properties of Superheavy Elements,'' \textit{Actinides Reviews}, vol. 1, pp. 433--485, 1971.
\bibitem{nist2024} A. Kramida et al., \textit{NIST Atomic Spectra Database} (ver. 5.11), National Institute of Standards and Technology, Gaithersburg, MD, 2024.
\bibitem{prism2026} D. A. Hayden and U. Ephemera, ``Project CODENAME PRISM: Discrete Phase-Space Relaxation on Commodity Silicon,'' Mage's Guild Psychonautics LLC Technical Report MGP-TR-2026-09, 2026.
\end{thebibliography}

\end{document}
